\subsection{A torsion bound and a theta-avoiding isotropic subgroup}

Let \(k\) be algebraically closed, let \((A,\Theta)\) be a
principally polarized abelian surface over \(k\), and suppose
\(\Theta\) is a smooth symmetric genus-two curve. Let
\(\ell\ge3\) be a prime different from \(\operatorname{char}k\).
The group \(A[\ell](k)\) is a four-dimensional
\(\mathbf F_\ell\)-space with the nondegenerate alternating
Weil pairing of the principal polarization; isotropic below
refers to this pairing.

\begin{lemma}[A direct bound on torsion points of a theta curve]
\label{lem:root-theta-torsion-delta-bound}
\status{Candidate; complete direct proof supplied.}
Put \(t=\#(\Theta(k)\cap A[\ell](k))\). Then
\[
 t(t-1)\le2\ell^4-2,\qquad
 t\le T_\ell:={\left\lfloor
             \frac{1+\sqrt{8\ell^4-7}}2\right\rfloor}
       <2\ell^2-1.
\]
\end{lemma}
\begin{proof}
Let \(D\subset A\) be the reduced image of \(\Theta\)
under the multiplication morphism \([\ell]:A\to A\).
No nonzero translation stabilizes \(\Theta\): such a
translation would stabilize \(\mathcal O_A(\Theta)\),
whereas its polarization homomorphism has trivial kernel.
For each nonzero \(h\in A[\ell](k)\), the two distinct
irreducible curves \(\Theta\) and \(\Theta+h\) therefore
have finite intersection. Outside the union of these finitely
many finite intersections, \([\ell]|_\Theta\) is injective.
It is finite and separable, since \([\ell]\) is etale.
It is consequently birational onto \(D\), and the smooth
curve \(\Theta\) is the normalization of \(D\).

In numerical divisor classes, \([\ell]^*\Theta=\ell^2\Theta\)
and \([\ell]_*[\ell]^*\Theta=\ell^4\Theta\). The
generic degree just proved is one, so
\(D\equiv[\ell]_*\Theta=\ell^2\Theta\).
Adjunction on the abelian surface, with trivial canonical
bundle and \(\Theta^2=2\), gives
\[
 p_a(D)=\ell^4+1,\qquad
 \sum_{x\in D}\delta_x(D)=p_a(D)-g(\Theta)=\ell^4-1.
\]
At the origin the \(t\) torsion preimages give exactly
\(t\) distinct smooth formal branches: the ambient
multiplication morphism is etale at every preimage.
For completeness, a reduced plane-curve germ with smooth
branches \(f_1,\ldots,f_t\) has
\(\delta=\sum_{i<j}I(f_i,f_j)\).
This follows by induction from the exact sequence
\[
0\longrightarrow R/(fg)\longrightarrow
 R/(f)\oplus R/(g)\longrightarrow R/(f,g)\longrightarrow0,
 \qquad R=k[[u,v]],
\]
and comparison with the direct sum of the branch
normalizations. Every pair of distinct branches through
the origin has intersection multiplicity at least one.
Hence \(\delta_0(D)\ge\binom t2\), proving the first
inequality and its quadratic solution. Finally
\[
(2\ell^2-1)(2\ell^2-2)-(2\ell^4-2)
 =2(\ell^2-1)(\ell^2-2)>0.
\]
Thus the integer \(t\) is strictly smaller than
\(2\ell^2-1\), as asserted. The argument also allows
\(t=0\), when the origin is not on \(D\).
\end{proof}

\begin{proposition}[An isotropic subgroup avoiding the theta curve]
\label{prop:root-theta-avoiding-Lagrangian}
\status{Candidate; complete direct proof supplied.}
There is a maximal isotropic subgroup
\(H\subset A[\ell](k)\) such that
\(H\cap\Theta(k)\subset\{0\}\). More quantitatively,
the number of such subgroups is at least
\[
(\ell+1)\bigl(\ell^2+1-\lfloor T_\ell/2\rfloor\bigr)
 \ \ge\ 2(\ell+1).
\]
\end{proposition}
\begin{proof}
Call a one-dimensional subspace bad if it contains a
nonzero point of \(\Theta\). Symmetry pairs each such
point with its distinct negative, since \(\ell\) is odd.
Different lines are disjoint away from zero, so the
number of bad lines is at most \(\lfloor t/2\rfloor\).
This bound remains valid whether or not zero lies on
\(\Theta\).

There are \((\ell^4-1)/(\ell-1)=(\ell+1)(\ell^2+1)\)
lines in the four-space. A fixed line lies in exactly
\(\ell+1\) Lagrangian planes, because these planes
correspond to the lines of its two-dimensional
symplectic perpendicular quotient. Each Lagrangian
plane itself contains \(\ell+1\) lines. Counting
line-plane incidences therefore shows that the total
number of Lagrangian planes is
\((\ell+1)(\ell^2+1)\). At most
\((\ell+1)\lfloor t/2\rfloor\) contain a bad line,
by the union bound. Subtracting gives the stated
lower bound. The preceding lemma gives
\(T_\ell\le2\ell^2-2\), so this bound is positive.
Every remaining plane has the required intersection.
Since \(\ell\) is invertible in \(k\), its constant
subgroup of geometric points is also an etale subgroup
scheme of \(A[\ell]\).
\end{proof}

\begin{remark}[The unproved Cartier bridge is separate]
\label{rem:root-theta-avoidance-scope}
\status{Candidate scope; no new monodromy theorem.}
The proposition does not assert that the image of the
theta curve under the quotient by \(H\) has only nodes:
it does not exclude a coset of \(H\) meeting the curve
in three or more points. Nor does it identify any
Schrodinger-basis coefficient of the Cartier tensor
with a theta value. Those two assertions remain
separate obligations. No full-monodromy case is
deduced here, and this auxiliary result does not
change a main gate.

The argument was explored after reading Pareschi's
original proof of Theorem C for genus two in
\emph{Torsion points on theta divisors and
semihomogeneous vector bundles}, arXiv:2007.12252v3,
printed pp.~9--10. That source works over the complex
numbers and invokes a Seshadri constant. The proof
above instead supplies the local delta calculation
directly over the stated field. No characteristic-two
extension of its semihomogeneous-bundle theorems is
assumed. No mathematical program was needed for the
displayed exact identities and incidence count.
\end{remark}
